\scanpage{112}
\deleted{Or \deleted{\unsure{$D$}} $D$ \inserted{[considéré comme sousschéma de $X$ p.ex. avec structure induite réduite]} est \uline{universellement ouvert sur $S$}, en vertu du th. de Chevalley [---] donc si $D'$ est son image inverse dans $S'$, on a aussi
\[
 D'_{s'}\subset(\overline{D'_{\eta'}})_{s'}.
\]
Or on a $D'_{\eta'}=\Theta'_{s'}$ ensemblistement, d'où $D'_{s'}\unsure{:|=}\Theta'_{s'}$ ensemblt, OK.}
\margin{appliqué aux composantes irréductibles de $D$}

\deleted{Donc il existe une \deleted{\unsure{enti}} section $\bar s$ de $\underline{\operatorname{Pic}}_{X/S}$ sur $S$ qui prolonge $s$, et qui correspond donc à un diviseur \deleted{$D'$} de Cartier positif $\Theta$ sur $X$, transversal aux fibres. \deleted{\unsure{comme} $\Theta_\eta=$} \inserted{car $\Theta_\eta=D_\eta$ par construction.} Considérons le cycle divisoriel $D-\Theta$, il ne contient aucune composante d'une fibre, et est nul sur $X_\eta$, donc est nul [lemme immédiat], d'où $D=\Theta$, cqfd.}

\deleted{Dans le cas général [$D$ pas néc. $\geq0$] si $f$ est projectif, il suffit d'ajouter à $D$ une \deleted{\unsure{multiple}} \inserted{section} d'un}
\margin{111}
